\begin{thebibliography}{99}
\bibitem{conch2023} Lu, M. Y. et al. A visual-language foundation model for
computational pathology. \textit{Nat. Med.} \textbf{30}, 863--874 (2024).
\bibitem{virchow2023} Vorontsov, E. et al. A foundation model for clinical-grade
computational pathology and rare cancers detection. \textit{Nat. Med.} \textbf{30},
2924--2935 (2024).
\bibitem{tonekaboni2019} Tonekaboni, S., Joshi, S., McCradden, M. D. \& Goldenberg,
A. What clinicians want: contextualizing explainable machine learning for clinical end
use. \textit{Proc. Mach. Learn. Res.} \textbf{106}, 359--380 (2019).
\bibitem{prinster2024} Prinster, D. et al. Care to explain? AI explanation types
differentially impact chest radiograph diagnostic performance and physician trust in AI.
\textit{Radiology} \textbf{313}, e233261 (2024).
\bibitem{kim2018} Kim, B. et al. Interpretability beyond feature attribution:
quantitative testing with concept activation vectors (TCAV). \textit{Proc. Mach. Learn.
Res.} \textbf{80}, 2668--2677 (2018).
\bibitem{koh2020} Koh, P. W. et al. Concept bottleneck models. \textit{Proc. Mach. Learn.
Res.} \textbf{119}, 5338--5348 (2020).
\bibitem{sauter2022} Sauter, D. et al. Validating automatic concept-based explanations
for AI-based digital histopathology. \textit{Sensors} \textbf{22}, 5346 (2022).
\bibitem{lou2026} Lou, P. et al. A human-in-the-loop explanation framework for
morphologically transparent AI predictions from whole-slide images. \textit{npj Digit.
Med.} \textbf{9}, 579 (2026).
\bibitem{han2026} Han, T. et al. CLEAR: an auditable foundation model for radiology
grounded in clinical concepts. \textit{Nat. Biomed. Eng.} (2026).
\bibitem{panda2022} Bulten, W. et al. Artificial intelligence for diagnosis and
Gleason grading of prostate cancer: the PANDA challenge. \textit{Nat. Med.}
\textbf{28}, 154--163 (2022).
\bibitem{nadt2021} Wilkinson, S. et al. Nascent prostate cancer heterogeneity drives
evolution and resistance to intense hormonal therapy [Data set]. The Cancer Imaging
Archive (2020).
\bibitem{dawood2026} Dawood, M. et al. Confounding factors and biases abound when
predicting molecular biomarkers from histological images. \textit{Nat. Biomed. Eng.}
(2026).
\bibitem{neidlinger2025} Neidlinger, P. et al. Benchmarking foundation models as
feature extractors for weakly supervised computational pathology. \textit{Nat. Biomed.
Eng.} \textbf{10}, 1113--1123 (2026).
\bibitem{pathorob2026} K\"omen, J. et al. Towards robust foundation models for digital
pathology. \textit{Nat. Commun.} \textbf{17}, 5218 (2026).
\bibitem{leopard2026} Faryna, K. et al. Predicting biochemical recurrence from
prostatectomy slides---the LEOPARD challenge. MIDL Short Papers Poster (2026).
\bibitem{grisi2026} Grisi, C. et al. Deep learning from routine histology improves risk
stratification for biochemical recurrence in prostate cancer. arXiv:2603.14187 (2026).
\bibitem{precise2026} Calapaqu\'i Ter\'an, A. K. et al. PRECISE---PRostate
Expert-annotated Contiguous IHC--H\&E Serial sEctions [Dataset]. Zenodo (2026).
\bibitem{tcgaprad2015} Cancer Genome Atlas Research Network. The molecular taxonomy of
primary prostate cancer. \textit{Cell} \textbf{163}, 1011--1025 (2015).
\bibitem{tcgacdr2018} Liu, J. et al. An integrated TCGA pan-cancer clinical data
resource to drive high-quality survival outcome analytics. \textit{Cell} \textbf{173},
400--416.e11 (2018).
\end{thebibliography}
